\documentclass[aspectratio=169]{beamer}
\mode<presentation> {
\usetheme{Madrid}
}

\usepackage{graphicx}
\usepackage{booktabs}
\usepackage{multicol}
\usepackage{multirow}
\usepackage{longtable}
\usepackage[utf8]{inputenc}
\usepackage{enumerate}
\usepackage{color,url}
\usepackage{tikz}
\usetikzlibrary{patterns}
\usepackage{calc}
\usepackage{subfig}
\usepackage{pgfplots}
\pgfplotsset{compat=1.9}
\usepackage{pgf-pie}
\usepackage{algorithm}
\usepackage{algpseudocode}
\usepackage{listings}
\lstset{language=Python,basicstyle=\ttfamily,keywordstyle=\color{blue}}
\usepackage{caption}

\AtBeginSection[]{
  \begin{frame}
  \vfill
  \centering
  \begin{beamercolorbox}[sep=8pt,center,shadow=true,rounded=true]{title}
    \usebeamerfont{title}\insertsectionhead\par%
  \end{beamercolorbox}
  \vfill
  \end{frame}
}

%----------------------------------------------------------------------------------------
%	TITLE PAGE
%----------------------------------------------------------------------------------------
\title{UnBlabla Car}
\subtitle{Mobilidade e Networking na UnB}

\author[Equipe]{Antonio Paulo Abreu Marques - 19008449 \newline Lucas Tavares Drumond - 231011650 \newline Heitor Fernandes Estevam - 231011453 \newline Luiz Felipe Bessa Santos - 231011687 \newline Gabriel Costa Santana - 190087391}

\institute[UnB]{Sistemas de Informação - CIC0101 \\ 
    \medskip
    {\includegraphics[scale=.25]{logo-unb}} 
    \medskip
    \date{\today}
}

\begin{document}

\begin{frame}
\titlepage
\end{frame}

\begin{frame}
\frametitle{Visão Estratégica}
\tableofcontents
\end{frame}

%------------------------------------------------
\section{Produto}
%------------------------------------------------

\begin{frame}
\frametitle{Problema}
Milhares de estudantes da UnB realizam trajetos diários semelhantes, muitas vezes com assentos vazios em seus veículos ou enfrentando dificuldades no transporte público como superlotação e atrasos.
\begin{itemize}
    \item \textbf{Desafio:} Como conectar motoristas e passageiros da comunidade acadêmica de forma segura e eficiente?
    \item \textbf{Solução (UnBlabla Car):} Uma plataforma peer-to-peer focada exclusivamente na comunidade da UnB.
    \item \textbf{Principais Aspectos:} 
    \begin{itemize}
        \item Otimização de tempo e recursos financeiros.
        \item Fortalecimento do networking acadêmico.
        \item Promoção de maior união entre os estudantes de forma orgânica
    \end{itemize}
\end{itemize}
\end{frame}

\begin{frame}
\frametitle{Escopo do Sistema}
\begin{columns}[c]
\column{.45\textwidth}
\textbf{O que o sistema É:}
\begin{itemize}
    \item Um ambiente fechado e inclusivo projetado para alunos da universidade.
    \item Um mecanismo de divisão de custos (gasolina) para incentivar a congregação acadêmica.
    \item Uma plataforma com foco mobile, priorizando a praticidade da mobilidade urbana.
\end{itemize}

\column{.5\textwidth}
\textbf{O que o sistema NÃO É:}
\begin{itemize}
    \item Não é um serviço de transporte comercial como Uber ou 99POP.
    \item Não visa a obtenção de lucro pelos motoristas ou pela plataforma.
    \item Não é aberto ao público geral, o uso é estrito aos alunos.
\end{itemize}
\end{columns}
\end{frame}

%------------------------------------------------
\section{Arquitetura e Modelagem de Dados}
%------------------------------------------------

\begin{frame}
\frametitle{Lógica de Banco de Dados}
\begin{block}{Estrutura Simples, Resolução Complexa}
A arquitetura foi desenhada para garantir performance e integridade:
\begin{itemize}
    \item \textbf{Disponibilidade e Desempenho:} Buscas rápidas e acessíveis durante picos (manhã e tarde).
    \item \textbf{Privacidade por Design:} O sistema protege o histórico de rotas e dados sensíveis dos usuários.
\end{itemize}
\end{block}
\vspace{0.5cm}
\end{frame}

% Utilização da tag [plain] para maximizar a área da imagem
\begin{frame}[plain]
\begin{figure}[htb!]
\centering
\includegraphics[width=\textwidth, height=0.95\textheight, keepaspectratio]{modelo-conceitual.png}
\end{figure}
\end{frame}

\begin{frame}[plain]
\begin{figure}[htb!]
\centering
\includegraphics[width=\textwidth, height=0.95\textheight, keepaspectratio]{modelo-fisico.png}
\end{figure}
\end{frame}

%------------------------------------------------
\section{Épicos e Histórias de Usuário}
%------------------------------------------------

\begin{frame}
\frametitle{Segurança e Pareamento}
\begin{block}{E1 — Acesso e Identidade}
    \begin{itemize}
        \item \textbf{Autenticação (H1):} Cadastro e acesso estritamente via e-mail institucional ativo.
        \item \textbf{Perfil (H2):} Cadastro de dados do estudante e do veículo.
    \end{itemize}
\end{block}

\begin{block}{E2 — Pareamento entre Usuários}
    \begin{itemize}
        \item \textbf{Rotas (H3/H4):} O motorista publica seu destino e os passageiros buscam rotas de saída próximas. O algoritmo deve sugerir um ponto de encontro sem grande desvio do caminho.
        \item \textbf{Cálculo Sem Atrito (H4):} O sistema sugere o valor de contribuição baseado na rota, eliminando o constrangimento.
    \end{itemize}
\end{block}
\end{frame}

\begin{frame}
\frametitle{Operação e Reputação}
\begin{block}{E3 — Logística da Viagem}
    \begin{itemize}
        \item \textbf{Confirmação (H5/H7):} Ciclo de vida da solicitação (pedido, aceite, conclusão/cancelamento).
        \item \textbf{Comunicação Segura (H6):} Chat provisório apenas para o embarque, protegendo o número de telefone pessoal dos usuários.
    \end{itemize}
\end{block}

\begin{block}{E4 — Confiança e Futuro}
    \begin{itemize}
        \item \textbf{Feedback (H8/H9):} Avaliação via estrelas (1 a 5) e comentários, além de sistema de denúncia e bloqueio.
        \item \textbf{Gamificação (Visão Futura):} Conquista de troféus ou selos e cálculo de economia de CO2 para incentivar o uso.
    \end{itemize}
\end{block}
\end{frame}

%------------------------------------------------
\begin{frame}
\Large{\centerline{Obrigado!}}
\Large{\centerline{Dúvidas?}}

\vspace{1cm}

\centering{
    \color{blue}Sistemas de Informação - UnB
}
\end{frame}
%------------------------------------------------
\end{document}
